\subsubsection{Subproceso 2: Selección y Registro de Cursos}



Para el levantamiento de información del Subproceso 2 se emplearon de forma complementaria dos técnicas: la \textbf{revisión documental} y la \textbf{observación directa} del sistema.

La revisión documental permitió identificar el marco normativo que regula la selección de cursos: la Resolución Rectoral RR-0570-2022 (tabla de límites de créditos), el Texto Único de Procedimientos Administrativos (TUPA) de la UNI y el Estatuto y Reglamento de Matrícula vigentes. Esta técnica fue adecuada porque las reglas de negocio del subproceso se encuentran formalizadas en instrumentos normativos de acceso institucional, lo que garantiza la precisión y trazabilidad de la información obtenida.

La observación directa consistió en el seguimiento del flujo de interacción del estudiante con la Intranet DIRCE durante el periodo de matrícula, registrando las pantallas, mensajes del sistema, restricciones activas y los pasos efectivos de la selección. Esta técnica fue pertinente porque el subproceso opera casi en su totalidad de forma mediada por la plataforma tecnológica, y la observación directa reveló comportamientos del sistema y puntos de fricción que los documentos normativos no recogen explícitamente (como la degradación del rendimiento en horas pico o la desaparición de asignaturas en la interfaz por omisiones de programación).


